\section{Conclusions}
\label{sec:conclusions}

Version 4.9.7 applies a frozen, pull-blind two-stage support protocol to
full-2016 inference.  Phase~1 tests lower edges from 28 through 33~MeV, plus a
34~MeV geometry control, using 25 paired conditional backgrounds.  No edge passes:
each has only 3 of 12 cell means below 0.75 in magnitude, only 1 of 4
background-only means below 0.75, and a worst mean pull between 2.281 and 2.696.
The 30 and 32~MeV lanes additionally contain one retained irreproducible refit each.
No support-specific full-data fitted amplitude, local $p_0$, or upper limit enters
this decision.

The failure is conditional on the generating construction.  The pre-existing 2016
10\% development sample/subset supplies partial observed-shape information, while
the full-2016 array enters before the decision only through one scalar normalization
over 26--210~MeV.  The hybrid mean also uses an explicitly waived broad-tail
continuation whose archived ROOT record has \texttt{fit\_ok=false}.  Its favorable
Pearson statistic, finite positive shape, in-bound parameters, and bitwise toy
reproduction do not turn it into a qualified physical truth or coverage generator.

A uniform post-decision fit of the unfluctuated analytic mean shows the same
mass-dependent mismatch at every edge: roughly $+2\sigma_A$ at 49~MeV,
$-2.6\sigma_A$ at 54~MeV, and $+3.5\sigma_A$ at 59~MeV.  This confirms that the
observed Phase-1 pattern is not merely a Poisson fluctuation.

\begin{samepage}
The predeclared stop action therefore authorizes no downstream production.  No
provisional support, Phase~2 confirmation, 65~MeV holdout, support-specific full-2016
observed scan, or reviewed 2016 state ledger is present.

\noindent\parbox{\linewidth}{The v4.9.7 combined upper limit and 100-toy combined
band are also absent by construction.}  Choosing the
numerically smallest failing edge would be a retrospective rule change and is not
done.
\end{samepage}

The separate 2021 robustness audit resolves the apparent ``disappearance'' of the
earlier 65~MeV signal-like feature more precisely.  The observed counts are unchanged,
but the fitted narrow-signal coefficient falls by 39.0\% and its local asymptotic
$Z$ changes from 4.252 to 2.403 when the GP-support prescription changes.  The support
change shifts the coarse-rebin phase and migrates the preferred GP kernel state; the
new continuum is higher through the masked region.  Optimizer repair and signal
normalization do not explain the change.  The defensible conclusion is that the
narrow-signal attribution is not robust to the support prescription---not that the
feature has been proved background or that signal has been excluded.

The v4.2 three-campaign combination and v4.9.5 2021-only scan remain preserved as
distinct historical states.  The unexecuted v4.9.7 observed and combined workflow
source is retained only to document the gate that prevented production.  A future
combined result requires a separately named, prospectively specified validation
protocol or an explicitly chosen historical-support fallback.  Neither is inferred
from this failed support scan.
